% Part: turing-machines
% Chapter: machines-computations
% Section: representing-tms

\documentclass[../../../include/open-logic-section]{subfiles}

\begin{document}

\olfileid{tur}{mac}{rep}
\olsection{ٹیورنگ مشیناں دی نمائندگی}

\begin{explain}
ٹیورنگ مشیناں نوں \emph{حالتی خاکیاں}
(state diagrams) نال تصویری شکل وچ
دکھایا جا سکدا اے۔ ایہناں خاکیاں
وچ حالتاں دے خانے تیراں نال جڑے
ہوندے نیں۔ جیویں توقع اے، ہر حالتی خانہ مشین دی
اک حالت دکھاندا اے۔ ہر تیر اوہ
حکم دکھاندا اے جیڑا اوس حالت توں
چلایا جا سکدا اے؛ حکم دی تفصیل
متعلقہ تیر دے اُتے یا تھلّے
لکھی ہوندی اے۔ ہِٹھلی مشین ویکھو:
ایس دیاں صرف دو اندرونیاں
حالتاں $q_0$ تے $q_1$ نیں،
تے اک حکم اے:
\[
\begin{tikzpicture}[->,>=stealth',shorten >=1pt,auto,node distance=2.8cm,
                    semithick]
  \tikzstyle{every state}=[fill=none,draw=black,text=black]

  \node[initial,state]   (A)              {$q_0$};
  \node[state]   (B) [right of=A] {$q_1$};

  \path (A) edge  node {\TMtrans{\TMblank}{\TMstroke}{\TMright}} (B);
\end{tikzpicture}
\]
یاد کرو کہ ٹیورنگ مشین کول
پڑھن-لکھن والا سر تے اک پٹی ہوندی اے
جیہدے اُتے انپٹ لکھیا ہوندا اے۔ ایس
حکم نوں ایداں پڑھو:
\emph{جے اک~$\TMblank$
حالت $q_0$ وچ پڑھیا جاوے،
تاں اک~$\TMstroke$ لکھو،
سر نوں سجّے ول سرکاؤ، تے حالت
$q_1$ وچ چلے جاؤ}۔
ایہ اوہو گل اے کہ
حالتی تبدیلی دا فنکشن
$\tuple{q_0, \TMblank}$
نوں $\tuple{q_1, \TMstroke,
\TMright}$ ول بھیجدا اے۔
\end{explain}

\begin{ex}
\emph{جفت مشین}: ہِٹھلی ٹیورنگ
مشین اُدوں تے صرف اُدوں
رُکدی اے جدوں پٹی اُتے
$\TMstroke$ دی گنتی جفت
ہووے (ایس مفروضے ہِٹھ کہ
سارے $\TMstroke$ پہلے
$\TMblank$ توں پہلاں آندے
نیں)۔
\[
\begin{tikzpicture}[->,>=stealth',shorten >=1pt,auto,node distance=2.8cm,
                    semithick]
  \tikzstyle{every state}=[fill=none,draw=black,text=black]

  \node[initial,state]         (A)              {$q_0$};
  \node[state]         (B) [right of=A] {$q_1$};

  \path (A) edge [bend left] node {\TMtrans{\TMstroke}{\TMstroke}{\TMright}} (B)
        (B) edge [loop above] node {\TMtrans{\TMblank}{\TMblank}{\TMright}} (B)
            edge [bend left] node {\TMtrans{\TMstroke}{\TMstroke}{\TMright}} (A);
\end{tikzpicture}
\]

ایہ حالتی خاکہ ہِٹھلے
حالتی تبدیلی دے فنکشن
دے مطابق اے:
\begin{align*}
\delta(q_0, \TMstroke) & = \tuple{q_1, \TMstroke, \TMright},\\
\delta(q_1, \TMstroke) & = \tuple{q_0, \TMstroke, \TMright},\\
\delta(q_1, \TMblank)  & = \tuple{q_1, \TMblank, \TMright}
\end{align*}
\end{ex}

\begin{explain}
اُتلی مشین صرف اُدوں
رُکدی اے جدوں انپٹ وچ
لکیراں دی گنتی جفت ہووے۔
ہور صورت وچ مشین نظری طور اُتے
ہمیشہ چلدی رہندی اے۔ کسے وی
مشین تے انپٹ لئی مشین
دیاں \emph{مکمل حالتاں}
(configurations) دا سلسلہ
ویکھ کے آؤٹ پٹ دا پتا
لایا جا سکدا اے۔ مکمل
حالت دی رسمی تعریف
اسی اگے دیواں گے۔
فی الحال ایہناں نوں چلن دے
وکھو وکھرے ویلیاں اُتے مشین دی حالت
دکھان والے خاکیاں دا
سلسلہ سمجھ لو۔ ہر مکمل
حالت پٹی دا متن، مشین
دی حالت تے پڑھن-لکھن
والے سر دی جگہ دکھاندی
اے۔

آؤ جفت مشین نوں چار
$\TMstroke$ دا انپٹ دے
کے اوس دیاں مکمل حالتاں
دا سلسلہ ویکھئیے۔ ایس
صورت وچ مشین دے رُکّن
دی توقع اے۔ فیر اسی
مشین نوں تِن $\TMstroke$
دے انپٹ اُتے چلاواں گے۔
اوس صورت وچ اوہ ہمیشہ چلدی
رہے گی۔

% PNB-CMP-026: source says “state one” although the initial index is zero.
مشین حالت~$q_0$ وچ
شروع ہوندی اے تے
سب توں کھبے پاسے والا~$\TMstroke$
پڑھدی اے۔ مشین دی
ابتدائی مکمل حالت
ایداں دکھا سکدے آں:
\[
\TMendtape \TMstroke_0 \TMstroke \TMstroke \TMstroke \TMblank \ldots
\]
اُتلی مکمل حالت
سمجھنی سوکھی اے۔
جیویں ویکھ سکدے آں،
مشین صفر نمبر والی
ابتدائی حالت وچ
سب توں کھبے پاسے والا~$\TMstroke$
پڑھدی اے۔ پہلی~$\TMstroke$
دے تھلّے حالتی ناں
دا ذیلی ہندسہ ایس
گل نوں دکھاندا اے۔
ایتھے لاگو ہون
والا حکم
$\delta(q_0, \TMstroke) = \tuple{q_1, \TMstroke, \TMright}$
اے؛ ایس لئی مشین دا سر
پٹی اُتے سجّے ول سرکدا اے تے مشین
حالت~$q_1$ وچ
چلی جاندی اے۔
\[
\TMendtape \TMstroke \TMstroke_1 \TMstroke \TMstroke \TMblank \ldots
\]
ہن مشین حالت~$q_1$
وچ اک~$\TMstroke$
پڑھ رہی اے؛ ایس
لئی اسی حکم
$\delta(q_1, \TMstroke) = \tuple{q_0,
  \TMstroke, \TMright}$
دی پیروی کرنی اے۔
ایس توں ہِٹھلی
مکمل حالت ملدی اے:
\[
\TMendtape \TMstroke \TMstroke \TMstroke_0 \TMstroke \TMblank \ldots
\]
مشین اگے چلدی اے
تاں اوہو قاعدے
اوہو ترتیب وچ
فیر لگدے نیں؛
اگلیاں دو مکمل
حالتاں ایہ نیں:
\[
\TMendtape \TMstroke \TMstroke \TMstroke \TMstroke_1 \TMblank \ldots
\]
\[
\TMendtape \TMstroke \TMstroke \TMstroke \TMstroke \TMblank_0 \ldots
\]
ہن مشین حالت~$q_0$
وچ اک~$\TMblank$
پڑھ رہی اے۔
حالتی تبدیلی دے
خاکے توں صاف اے
کہ ایتھے کوئی
حکم موجود نہیں؛
ایس لئی مشین
رُک گئی اے۔
ایہدا مطلب اے
کہ انپٹ قبول
ہو گیا اے۔

ہن فرض کرو اسی
مشین نوں تِن
$\TMstroke$ دا انپٹ
دے کے شروع کریے۔
پہلیاں کجھ مکمل
حالتاں پچھلیاں
ورگیاں نیں: اوہو
حکم چلدے نیں،
صرف پٹی دا انپٹ
تھوڑا وکھرا اے:
\[
\TMendtape \TMstroke_0 \TMstroke \TMstroke \TMblank \ldots
\]
\[
\TMendtape \TMstroke \TMstroke_1 \TMstroke \TMblank \ldots
\]
\[
\TMendtape \TMstroke \TMstroke \TMstroke_0 \TMblank \ldots
\]
\[
\TMendtape \TMstroke \TMstroke \TMstroke \TMblank_1 \ldots
\]
ہن مشین سارے
$\TMstroke$ لنگھ
چکی اے تے
اک~$\TMblank$
حالت~$q_1$ وچ
پڑھ رہی اے۔
خاکے مطابق
ایس لئی حکم
$\delta(q_1, \TMblank) =\tuple{q_1,
\TMblank, \TMright}$
موجود اے۔ چونکہ
پٹی سجّے پاسے
ول بے انتہا
$\TMblank$ نال
بھری اے، مشین
ایہو حکم
\emph{ہمیشہ}
چلاندی رہے گی،
حالت~$q_1$ وچ رہندی اے،
تے اوہدا سر ہور ہور سجّے ول سرکدا رہندا اے۔
اوہ کدی
نہیں رُکے گی
تے انپٹ قبول
نہیں کرے گی۔
\end{explain}

\begin{explain}
ایہ یاد رکھنا ضروری
اے کہ کجھ مشیناں رُکدیاں ای نہیں۔ جے
رُکّن دا مطلب ایہ اے
کہ مشین کول چلان لئی
حکم مک جان، تاں ہر
حالت وچ ہر نشان لئی
باہر جان والا تیر
رکھ کے اسی ایداں
دی مشین بنا سکدے آں
جیڑی کدی نہ رُکے۔
جفت مشین نوں وی
$\TMblank$ حالت~$q_0$
وچ پڑھݨ
لئی اک حکم پا کے
ہمیشہ چلن والی
بنا سکدے آں۔
\end{explain}

\begin{ex}
\[
\begin{tikzpicture}[->,>=stealth',shorten >=1pt,auto,node distance=2.8cm,
                    semithick]
  \tikzstyle{every state}=[fill=none,draw=black,text=black]

  \node[initial,state] (A)              {$q_0$};
  \node[state]         (B) [right of=A] {$q_1$};

  \path
  (A) edge [bend left] node {\TMtrans{\TMstroke}{\TMstroke}{\TMright}} (B)
      edge [loop above] node {\TMtrans{\TMblank}{\TMblank}{\TMright}} (A)
  (B) edge [loop above] node {\TMtrans{\TMblank}{\TMblank}{\TMright}} (B)
      edge [bend left] node {\TMtrans{\TMstroke}{\TMstroke}{\TMright}} (A);
\end{tikzpicture}
\]
\end{ex}

\begin{explain}
ٹیورنگ مشیناں دی نمائندگی
مشینی جدولاں (machine tables) نال وی ہو سکدی اے۔
ایہناں وچ پٹی دے
حرفاں دا مجموعہ
$x$-محور اُتے
تے مشین دیاں
حالتاں دا سیٹ
$y$-محور اُتے
دکھایا جاندا اے۔
جدول وچ ہر حالت دی قطار تے ہر نشان دے کالم
دے ملاپ والے خانے وچ حکم دا باقی حصہ،
یعنی نویں حالت،
نواں نشان تے
حرکت دی سمت،
لکھی ہوندی اے۔
مشینی جدول توں
ایہ لبھنا سوکھا
اے کہ مشین کس
حالت تے کس نشان
اُتے رُکدی اے۔
جدول وچ جتھے
خانہ خالی ہووے،
اوتھے مشین رُک
سکدی اے۔ حالتی
خاکیاں تے حکم دی
فہرست وچ رُکّن
والیاں جگہاں فوراً
ظاہر نہیں ہوندیاں؛
جدول دے خالی
خانے ویکھ کے
اوہ جلدی لبھ
جاندیاں نیں۔
\end{explain}

\begin{ex}
جفت مشین دا جدول ایہ اے:
\[
\centering
\begin{tabular}{lllll}
\cline{2-4}
\multicolumn{1}{l|}{}      & \multicolumn{1}{c|}{$\TMblank$}                
& \multicolumn{1}{c|}{$\TMstroke$}     & \multicolumn{1}{c|}{$\TMendtape$}   &  \\ \cline{1-4}
\multicolumn{1}{|l|}{$q_0$} & \multicolumn{1}{l|}{}                          
& \multicolumn{1}{l|}{\TMtrans{\TMstroke}{q_1}{\TMright}} & 
\multicolumn{1}{l|}{\phantom{\TMtrans{\TMstroke}{q_1}{\TMright}}} &  \\ \cline{1-4}
\multicolumn{1}{|l|}{$q_1$} & \multicolumn{1}{l|}{\TMtrans{\TMblank}{q_1}{\TMright}} 
& \multicolumn{1}{l|}{\TMtrans{\TMstroke}{q_0}{\TMright}} & 
\multicolumn{1}{l|}{\phantom{\TMtrans{\TMstroke}{q_1}{\TMright}}} &  \\ \cline{1-4}
\end{tabular}
\]
جیویں ویکھ سکدے آں،
مشین اک~$\TMblank$
حالت~$q_0$ وچ پڑھدے
ویلے رُکدی اے۔
\end{ex}

\begin{explain}
ہُن تک اسی صرف
اوہ مشیناں ویکھیاں
نیں جیہڑیاں انپٹ
پڑھدیاں تے قبول
کردیاں نیں۔ پر
ٹیورنگ مشیناں
پڑھ وی سکدیاں
تے لکھ وی سکدیاں
نیں۔ ایہدی اک مثال
(ہور بہت مثالاں وی
نیں) \emph{دوگنا کرن
والی مشین} (doubler)
اے۔ جے پٹی اُتے~$n$
$\TMstroke$ دا
اک ٹکڑا دے کے
ایہ مشین شروع
کریے، تاں ایہ
$2n$ $\TMstroke$
دا ٹکڑا آؤٹ پٹ
دیندی اے۔
\end{explain}

\begin{ex}
  \ollabel{ex:doubler} دوگنا
  کرن والی مشین بناؤن
  توں پہلاں مسئلہ
  حل کرن دی اک
  \emph{تدبیر} بنانی
  ضروری اے۔ جیویں
  اسی مشین نوں
  ہُن تک بیان کیتا
  اے، ایہ یاد نہیں
  رکھ سکدی کہ
  اوس نے کنے
  $\TMstroke$ پڑھے
  نیں؛ ایس لئی
  پٹی اُتے موجود
  سارے $\TMstroke$
  دا حساب رکھن
  دا کوئی طریقہ
  لبھنا پئے گا۔
  اک طریقہ ایہ
  اے کہ آؤٹ پٹ
  نوں انپٹ توں
  اک~$\TMblank$ نال
  وکھ کر دئیے۔
  مشین فیر انپٹ
  وچوں پہلا
  $\TMstroke$ مٹا
  دیوے، باقی انپٹ
  لنگھے، اک
  $\TMblank$ چھڈے
  تے دو نویں
  $\TMstroke$ لکھے۔
  فیر مشین واپس
  جا کے انپٹ وچ
  دوجا $\TMstroke$
  لبھے تے اوس
  نوں وی دوگنا
  کرے۔ انپٹ دے
  ہر اک $\TMstroke$
  لئی آؤٹ پٹ دے
  دو $\TMstroke$
  لکھے گی۔ چلن
  نال نال انپٹ
  مٹاندی جاوے تاں
  ایہ یقین رہندا
  اے کہ کوئی
  $\TMstroke$ نہ
  رہ جاوے تے نہ اوہنوں دو واری دوگنا
  کیتا جاوے۔ جدوں
  سارا انپٹ مٹ
  جاوے گا، پٹی
  اُتے $2n$
  $\TMstroke$ رہ
  جان گے۔ نتیجے
  والی ٹیورنگ
  مشین دا حالتی
  خاکہ \olref{fig:doubler}
  وچ دکھایا گیا اے۔
  \begin{figure}
\[
\begin{tikzpicture}[->,>=stealth',shorten >=1pt,auto,node distance=2.8cm,
                    semithick]
  \tikzstyle{every state}=[fill=none,draw=black,text=black]

  \node[initial,state] (1)              {$q_0$};
  \node[state]         (2) [right of=1] {$q_1$};
  \node[state]         (3) [right of=2] {$q_2$};
  \node[state]         (4) [below of=3] {$q_3$};
  \node[state]         (5) [left of=4]  {$q_4$};
  \node[state]         (6) [left of=5]  {$q_5$};

  \path (1) edge node {\TMtrans{\TMstroke}{\TMblank}{\TMright}} (2)
    (2) edge [loop above] node {\TMtrans{\TMstroke}{\TMstroke}{\TMright}} (2)
      edge node {\TMtrans{\TMblank}{\TMblank}{\TMright}} (3)
    (3) edge [loop above] node {\TMtrans{\TMstroke}{\TMstroke}{\TMright}} (3)
        edge node {\TMtrans{\TMblank}{\TMstroke}{\TMright}} (4)
    (4) edge [loop below] node {\TMtrans{\TMblank}{\TMstroke}{\TMleft}} (4)
        edge node {\TMtrans{\TMstroke}{\TMstroke}{\TMleft}} (5)
    (5) edge [loop below]  node {\TMtrans{\TMstroke}{\TMstroke}{\TMleft}} (5)
        edge              node {\TMtrans{\TMblank}{\TMblank}{\TMleft}} (6)
    (6) edge [loop below] node {\TMtrans{\TMstroke}{\TMstroke}{\TMleft}} (6)
        edge              node {\TMtrans{\TMblank}{\TMblank}{\TMright}} (1);
\end{tikzpicture}
\]
\caption{اک دوگنا کرن والی مشین}
\ollabel{fig:doubler}
\end{figure}
\end{ex}

\begin{prob}
کوئی وی انپٹ
چنو تے
\olref[tur][mac][rep]{ex:doubler}
وچ دتی دوگنا
کرن والی مشین
دیاں مکمل حالتاں
دا سلسلہ قدم بہ
قدم ویکھو۔
\end{prob}

\begin{prob}
حرفاں دے مجموعے
$\{\TMendtape,\TMblank, A, B\}$
والی ٹیورنگ مشین
بناؤ جیڑی $A$
تے $B$ دیاں اوہ
لڑیاں قبول کرے،
یعنی اوہناں اُتے
رُک جائے، جِنہاں
وچ $A$ دی گنتی
$B$ دی گنتی دے
برابر ہووے
\emph{تے} سارے
$A$، سارے $B$
توں پہلاں آندے
ہون۔ جِنہاں لڑیاں
وچ $A$ دی گنتی
$B$ دی گنتی
دے برابر نہ ہووے
یا سارے $A$،
سارے~$B$ توں
پہلاں نہ آندے
ہون، اوہناں نوں
رد کرے، یعنی
اوہناں اُتے نہ
رُکے۔ مثال لئی،
مشین $AABB$
تے $AAABBB$
قبول کرے، پر $AAB$
تے $AABBAABB$
دونوں نوں رد کرے۔
\end{prob}

\begin{prob}
حرفاں دے مجموعے
$\{\TMendtape,\TMblank, A, B\}$
والی ٹیورنگ مشین
بناؤ جیڑی کوئی
وی لڑی $\alpha$
انپٹ لے؛ ایہ
لڑی $A$ تے $B$
دی بنی ہووے۔
مشین اوس دی اک ہور نقل
اوہدے نال جوڑ کے
$\alpha\alpha$
دی صورت وچ
آؤٹ پٹ بناوے۔
مثال لئی، انپٹ
$ABBA$ توں
آؤٹ پٹ
$ABBAABBA$
بننا چاہیدا اے۔
\end{prob}

\begin{prob}
\emph{حرف وار ترتیب؟:}
حرفاں دے مجموعے
$\{\TMendtape,\TMblank, A, B\}$
والی ٹیورنگ مشین
بناؤ جیڑی
$A$ تے $B$
دا اک محدود
سلسلہ انپٹ
لے کے جانچے
کہ سارے $A$
سارے $B$
دے کھبے پاسے
نیں یا نہیں۔
مشین انپٹ
دی لڑی پٹی
اُتے رہن دیوے؛
جے لڑی «حرف
وار ترتیب وچ»
ہووے تاں رُکے،
نہ ہووے تاں
ہمیشہ چلدی
رہے۔
\end{prob}

\begin{prob}
\emph{حرف وار ترتیب
دین والی مشین:}
حرفاں دے مجموعے
$\{\TMendtape,\TMblank, A, B\}$
والی ٹیورنگ مشین
بناؤ جیڑی
$A$ تے $B$
دا اک محدود
سلسلہ انپٹ
لے کے اوس
دے حرف ایداں
فیر ترتیب دے
کہ سارے $A$
سارے~$B$ دے
کھبے پاسے
ہون۔ مثال لئی،
لڑی $BABAA$
لڑی $AAABB$
بن جائے، تے
لڑی $ABBABB$
لڑی $AABBBB$
بن جائے۔
\end{prob}

\end{document}
